\documentclass{article}

% ============================================================
% MA 213 In-Class Worksheet Template
% Student-facing worksheet for lecture activities.
% ============================================================

\usepackage[margin=1in]{geometry}
\usepackage{amsmath}
\usepackage{xcolor}
\usepackage{parskip}
\usepackage{array}
\usepackage{enumitem}
\usepackage{booktabs}

\newcommand{\coursenumber}{MA 213}
\newcommand{\weekplaceholder}{5}
\newcommand{\lectureplaceholder}{12}
\newcommand{\lecturetitle}{Using the Geometric Distribution: Think/Pair/Share}

\newcommand{\nameblank}{\rule{1.75in}{0.4pt}}
\newcommand{\idblank}{\rule{1.55in}{0.4pt}}
\newcommand{\shortblank}{\rule{1in}{0.4pt}}
\newcommand{\longblank}{\rule{0.93\linewidth}{0.4pt}}

\begin{document}

\begin{center}
  {\LARGE \coursenumber{} Week \weekplaceholder{}: Lecture \lectureplaceholder{}}\\
  {\large \lecturetitle{}}
\end{center}

\begin{enumerate}[label=\arabic*.,leftmargin=*,itemsep=.7em]
  \item \textbf{Your name:} \nameblank \hfill \textbf{BU ID:} \idblank
\end{enumerate}

\textbf{Big idea:} $X \sim \text{Geometric}(p)$ gives the trial number of the \emph{first} success in a sequence of independent trials, with $Pr(X=n) = (1-p)^{n-1}p$. Practice setting this up in a new context, then apply it back to the Milgram experiment from earlier in this lecture.

\section*{Step 1: Think (2 mins)}

\textbf{(a)} A website randomly displays one of 30 different motivational quotes each time you visit. Only one of these quotes is by Maya Angelou. You refresh the page repeatedly, each time seeing a randomly selected quote. What is the probability that the first Maya Angelou quote appears on your 5th refresh? Give your answer as a formula in terms of powers and fractions.
\vspace{1.6in}

\newpage
\section*{Step 2: Pair (2 mins)}

\begin{enumerate}[label=\arabic*.,leftmargin=*,itemsep=.7em, start=2]
  \item \textbf{Partner's name:} \nameblank \hfill \textbf{BU ID:} \idblank
\end{enumerate}

Compare your answer to (a) with your partner, then work through this together.

\textbf{(b)} Recall that in the Milgram experiment, about 65\% of participants obeyed the experimenter and administered the shock, so about 35\% refused. Suppose Milgram tests participants one at a time, and (based on this historical rate) each participant independently has a 35\% chance of being the first to refuse. What is the probability that the first refusal happens on the 3rd participant tested? Set up $X$ and compute $Pr(X=3)$.
\vspace{0.7in}

\textbf{(c)} What assumptions does modeling this with a Geometric distribution require (think about the conditions for a sequence of Bernoulli trials)? Are these assumptions realistic for real human participants tested one after another?
\vspace{0.5in}

\textbf{Our thoughts:}
\vfill

\end{document}
